\documentclass[12pt]{article}
\usepackage[margin=1in]{geometry}
\usepackage{setspace}
\usepackage{amsmath,amssymb}
\usepackage{graphicx}
\usepackage{enumitem}
\usepackage{booktabs}
\usepackage{hyperref}
\usepackage{fontspec}
\usepackage{polyglossia}
\setdefaultlanguage{english}
\setmainfont{TeX Gyre Termes}
\setsansfont{TeX Gyre Heros}
\setmonofont{TeX Gyre Cursor}

\title{Summary of the ISAC Channel and SINR Simulation Project}
\author{Generated from the MATLAB/Octave codebase}
\date{\today}

\begin{document}
\maketitle

\begin{abstract}
This project implements a simplified simulation framework for an integrated sensing and communication (ISAC) system. The code generates communication and sensing channels, creates transmitted signals, models interference and noise, and evaluates the signal-to-interference-plus-noise ratio (SINR) for the legitimate receiver (Bob). The structure is organized into configuration, channel generation, signal generation, receiving modules, and performance metrics.
\end{abstract}

\section{Overview}
The project is designed to model a multi-user downlink system in which a base station transmits signals to several legitimate users while also considering an eavesdropper and a radar-type sensing channel. In the current implementation, the simulation focuses on:

\begin{itemize}[leftmargin=*]
    \item generating random complex channel matrices,
    \item creating transmitted symbols,
    \item adding artificial noise and receiver noise,
    \item evaluating Bob's received SINR,
    \item providing a simple baseline simulation framework for further optimization.
\end{itemize}

The current code is therefore best understood as a baseline simulation and analysis template for an ISAC or secure beamforming study, rather than a complete optimization engine.

\section{System Configuration}
The configuration is stored in \texttt{config/sim	extunderscore config.m}. It defines the system sizes and simulation parameters, including:

\begin{itemize}[leftmargin=*]
    \item $N_t = 30$: number of transmit antennas,
    \item $N_r = 20$: number of radar receive antennas,
    \item $K = 6$: number of legitimate users (Bob users),
    \item $J = 6$: number of eavesdroppers (Eve users),
    \item $T = 800$: number of transmitted symbols per block.
\end{itemize}

It also defines the noise variances:
\begin{align*}
\sigma_C^2 &= 1 \quad \text{(Bob noise variance)} \\
\sigma_I^2 &= 1 \quad \text{(Eve noise variance)} \\
\sigma_S^2 &= 1 \quad \text{(sensing noise variance)}.
\end{align*}

The file also sets performance targets and derived quantities such as $\gamma$, $\rho$, and $\beta$ from dB values.

\section{Channel Generation}
The channel model is implemented in \texttt{channel/generate	extunderscore channel.m}. It creates three matrices:

\begin{itemize}[leftmargin=*]
    \item $H \in \mathbb{C}^{K \times N_t}$: channel from base station to Bob users,
    \item $F \in \mathbb{C}^{J \times N_t}$: channel from base station to Eve users,
    \item $G \in \mathbb{C}^{N_r \times N_t}$: sensing / radar-related channel.
\end{itemize}

Each matrix is generated using i.i.d. complex Gaussian entries normalized by $\sqrt{2}$, which yields a Rayleigh-like fading model commonly used in wireless channel simulation.

\section{Signal Generation}
The symbol stream is generated by \texttt{signal/generate	extunderscore symbols.m}. It creates a matrix $S \in \mathbb{C}^{K \times T}$ containing random complex symbols for each user.

The transmitted signal is constructed in \texttt{signal/transmit	extunderscore signal.m}. In this implementation, the code:

\begin{itemize}[leftmargin=*]
    \item initializes a zero transmitted signal matrix $X \in \mathbb{C}^{N_t \times T}$,
    \item loops over each user, forming a beamforming contribution,
    \item adds artificial noise through \texttt{generate	extunderscore artificial	extunderscore noise.m}.
\end{itemize}

In the current version, the beamformer is a simple baseline using the identity matrix for each user, so this step acts as a basic reference rather than a fully optimized beamforming design.

\section{Reception and Noise Modeling}
The received signal at Bob is modeled in \texttt{receiver/receive	extunderscore bob.m}:
\begin{equation}
Y_{\text{Bob}} = H X + V_{\text{Bob}},
\end{equation}
where $V_{\text{Bob}}$ is complex Gaussian noise with variance determined by $\sigma_C^2$.

Similarly, the eavesdropper observation is built in \texttt{receiver/receive	extunderscore eve.m}:
\begin{equation}
R_{\text{Eve}} = F X + V_{\text{Eve}}.
\end{equation}

These equations represent the wireless channel plus additive noise, which is the standard model used in secure communication and physical-layer security simulations.

\section{SINR Metrics}
The project includes two SINR functions related to Bob.

\subsection{Simulation-based SINR}
The file \texttt{metrics/SINR	extunderscore bob	extunderscore sim.m} computes a Monte Carlo-style approximation of the SINR by comparing the received signal power aligned with the transmitted symbol to the remaining power in the received signal.

The logic is:
\begin{itemize}[leftmargin=*]
    \item compute the received vector for user $k$,
    \item multiply it by the conjugate of the symbol,
    \item estimate useful signal power,
    \item divide by the residual interference-plus-noise power.
\end{itemize}

This yields a practical simulation-based estimate of the effective SINR.

\subsection{Theoretical SINR}
The file \texttt{metrics/SINR	extunderscore bob	extunderscore theory.m} implements a theoretical approximation based on matrix traces. It uses the channel covariance term $H_k^T H_k$ and evaluates the signal power versus interference from other users plus noise:
\begin{equation}
\text{SINR}_k \approx \frac{\mathrm{tr}(W_k H_k)}{\sum_{j \neq k} \mathrm{tr}(W_j H_k) + \sigma_C^2}.
\end{equation}

This is a simplified analytical expression used to benchmark the simulation against an approximate closed-form estimate.

\section{Main Simulation Flow}
The script \texttt{main.m} ties all components together. The flow is:

\begin{enumerate}[leftmargin=*]
    \item load configuration,
    \item generate channel matrices,
    \item generate user symbols,
    \item build beamformers and artificial noise,
    \item transmit the signal,
    \item receive at Bob and Eve,
    \item compute simulation and theoretical SINR,
    \item print summary statistics.
\end{enumerate}

The main script is therefore the entry point for the experiment and acts as the central orchestrator of the complete simulation pipeline.

\section{What the Code Accomplishes}
At a high level, the code implements the following:

\begin{itemize}[leftmargin=*]
    \item a stochastic wireless channel model,
    \item a multi-user communication signal model,
    \item artificial noise and receiver noise injection,
    \item Bob-side SINR estimation,
    \item a simple simulation framework suitable for secure ISAC or beamforming research.
\end{itemize}

\section{Important Practical Note}
The repository currently contains a clean simulation framework, but it is not yet a complete optimized secure beamforming solution. The beamforming step is intentionally kept simple, and the evaluation focuses on demonstrating the data flow and metric computation rather than solving a full optimization problem. In other words, the code is a strong baseline and a starting point for further research and algorithm development.

\section{Conclusion}
This project provides a compact MATLAB/Octave simulation pipeline for a multi-user communication system with sensing and eavesdropper channels. The main contribution is the end-to-end modeling of channel generation, signal transmission, reception, and SINR evaluation, which makes it useful as a foundation for further work on beamforming, secrecy, and ISAC optimization.

\end{document}
